% Data Availability Statement for P5 (Digital Discovery submission)
\documentclass[11pt,a4paper]{article}

\usepackage[margin=1in]{geometry}
\usepackage[T1]{fontenc}
\usepackage[utf8]{inputenc}
\usepackage{lmodern}
\usepackage{booktabs}
\usepackage[colorlinks=true,linkcolor=blue,citecolor=blue,urlcolor=blue]{hyperref}
\usepackage{enumitem}

\title{Data Availability Statement\\[0.5em]
\large Chemical distribution shift in antimalarial natural products:\\
fingerprint dominance and localized topological complementarity\\
in molecular representation learning}

\author{}
\date{}

\begin{document}

\maketitle

\section*{Data Availability Statement}

All data and computational resources supporting the findings of this study are publicly available through Zenodo, ensuring complete reproducibility and compliance with FAIR principles (Findable, Accessible, Interoperable, Reusable).

\subsection*{Primary Data Repository}

\textbf{Zenodo Deposit:} All computational artifacts, including molecular datasets, trained models, analysis scripts, and complete results, are archived at:

\begin{center}
\url{https://doi.org/10.5281/zenodo.22787438}
\end{center}

\noindent \textbf{License:} Creative Commons Attribution 4.0 International (CC BY 4.0)

\noindent The Zenodo deposit includes:

\begin{itemize}[leftmargin=*,itemsep=0.5em]
\item \textbf{Molecular panel} (\texttt{data/p5\_panel\_19836.csv}): Complete SMILES strings, canonical identifiers, and binary activity labels for all 19,836 molecules in the benchmark panel, derived from 396 African natural products (ANPDB/AfroDb) and 454 synthetic antimalarials.

\item \textbf{Chemical splits} (\texttt{data/p5\_splits/}): Frozen train/test indices for all evaluation partitions (random, canonical scaffold, three novel scaffold variants, Butina cluster $C=0.55$) across 5 folds and 5 random seeds, enabling exact replication of all reported metrics.

\item \textbf{Pre-computed graph representations} (\texttt{data/p5\_graph\_cache.pkl}): PyTorch Geometric molecular graphs with node features (atom type, degree, hybridization, aromaticity) and edge indices for all 19,836 molecules, eliminating the need to reconstruct graph topology.

\item \textbf{Benchmark results} (\texttt{results/}): Complete performance records including:
\begin{itemize}
\item Per-fold predictions for all 25 fold--seed combinations per model (ECFP4-RF, GIN, GIN-TFP, GIN-TNE, ChemBERTa) under random and scaffold splits
\item Aggregated statistics (\texttt{p5\_replication\_stats.csv}) with per-seed means, population standard deviations, and Benjamini--Hochberg corrected $p$-values
\item External validation on molecule-disjoint ChEMBL panel (22,267 compounds, \texttt{p5\_public\_chembl\_malaria\_disjoint.csv})
\end{itemize}

\item \textbf{Extended robustness campaign} (\texttt{results/extended\_campaign\_20260825/}): Secondary analyses comprising 625 fold--seed records across 25 configurations, including descriptor permutation tests, salience stability matrices, and novel scaffold partition families.

\item \textbf{Distribution-shift analyses} (\texttt{results/nn\_tanimoto\_deciles\_20260829/}): Nearest-neighbour Tanimoto similarity stratification with per-decile AUC values for all models, demonstrating distance-resolved performance gradients across the chemical distribution shift spectrum.

\item \textbf{Calibration records} (\texttt{results/calibration\_20260827/}): Expected calibration error (ECE), maximum calibration error (MCE), and Brier scores for all 30 configurations (25 GNN-family + 5 ChemBERTa) under pooled test predictions, plus isotonic recalibration parameters.

\item \textbf{Butina scaffold-cluster benchmark} (\texttt{results/butina\_cluster\_20260829/}): Stricter chemical-extrapolation test with per-fold predictions, cluster assignments, and summary statistics for all five arms.

\item \textbf{Training scripts} (\texttt{scripts/}): Python implementations of all model architectures (GIN, GIN-TFP, GIN-TNE, ChemBERTa), baseline methods (ECFP4-RF, kNN, logistic regression), and analysis pipelines with command-line interfaces for reproducibility.

\item \textbf{Computational environment} (\texttt{scripts/environment.yml}): Complete conda environment specification (Python 3.11, PyTorch 2.13.0, PyTorch Geometric 2.8.0, Transformers 5.14.1, RDKit 2025.03.6, scikit-learn 1.5.1) with pinned dependency versions.

\item \textbf{Documentation} (\texttt{documentation/}): Data Analysis Report (\texttt{P5\_DATA\_ANALYSIS\_REPORT.md}) documenting complete computational provenance, including job identifiers, execution timestamps, quality-control gates, and evidence status labels for all reported results.

\item \textbf{Integrity verification} (\texttt{sha256sums.txt}): SHA256 checksums for all files in the deposit, enabling cryptographic verification of data integrity.
\end{itemize}

%\subsection*{Version Control Repository}
%
%The complete project repository, including manuscript source files, analysis scripts, and developmental history, is maintained on GitHub:
%
%\begin{center}
%\url{https://github.com/NanaEngo/Malaria_codesV2}
%\end{center}
%
%\noindent Specific files relevant to this study are located in:
%
%\begin{itemize}[leftmargin=*]
%\item \texttt{Project5\_GNN\_Transformer\_DrugDiscovery\_V2609/}
%\end{itemize}
%
%\noindent The GitHub repository provides:
%\begin{itemize}[leftmargin=*,itemsep=0.3em]
%\item Full version history with timestamped commits documenting all computational decisions
%\item Issue tracking and developmental annotations
%\item Integration with continuous workflow documentation
%\item Live synchronization with working computational environment
%\end{itemize}

\subsection*{Upstream Data Sources}

The molecular panel derives from:

\begin{enumerate}[leftmargin=*,itemsep=0.5em]
\item \textbf{African Natural Products Database (ANPDB):} 396 African natural product structures from \url{https://african-compounds.org/anpdb/} (accessed January 2026). ANPDB is a freely accessible repository of natural products from African medicinal plants, providing SMILES strings, structural descriptors, and ethnopharmacological annotations.

\item \textbf{AfroDb:} Supplementary African natural product structures from \url{https://africancompounds.org/} (accessed January 2026), providing additional coverage of traditional African medicinal chemistry scaffolds.

\item \textbf{ChEMBL antimalarial compounds:} 454 synthetic antimalarial molecules with reported \emph{Plasmodium falciparum} activity (ChEMBL release 33, accessed February 2026), retrieved via ChEMBL web services (\url{https://www.ebi.ac.uk/chembl/}).

\item \textbf{External validation panel:} 22,447 unique compounds from ChEMBL with reported antimalarial activity against \emph{P. falciparum}, filtered to 22,267 molecule-disjoint structures after removing 180 exact canonical-SMILES overlaps with the training panel (\texttt{p5\_public\_chembl\_malaria\_disjoint.csv}).
\end{enumerate}

\noindent All molecular structures underwent canonical SMILES standardization, tautomer normalization, and stereochemistry retention using RDKit (version 2025.03.6) with documented processing pipelines archived in the Zenodo deposit.

\subsection*{Model Checkpoints and Weights}

Due to file size constraints (trained model checkpoints: $\sim$8 GB), full model weights are available upon request from the corresponding author (\texttt{myke-vital.sao@facsciences-uy1.cm}). The Zenodo deposit includes:

\begin{itemize}[leftmargin=*,itemsep=0.3em]
\item Training hyperparameters and architecture specifications for exact replication
\item Random seed values ensuring deterministic weight initialization
\item Per-fold training logs with epoch-level loss and validation metrics
\item Final per-fold test predictions (probability scores) enabling metric recomputation without retraining
\end{itemize}

\noindent Alternatively, researchers may retrain all models using the provided scripts and frozen chemical splits, which will reproduce the reported results within expected stochastic variation (mean AUC differences $<0.002$ across independent training runs with identical seeds).

\subsection*{Software Dependencies}

All computational analyses were conducted in the \texttt{malaria\_md} conda environment with the following core dependencies:

\begin{itemize}[leftmargin=*,itemsep=0.2em]
\item Python 3.11.5
\item PyTorch 2.13.0 (CUDA 12.1)
\item PyTorch Geometric 2.8.0
\item Transformers 5.14.1 (Hugging Face)
\item RDKit 2025.03.6
\item scikit-learn 1.5.1
\item NumPy 1.26.4
\item Pandas 2.2.2
\item SciPy 1.13.0
\end{itemize}

\noindent The complete environment specification with all transitive dependencies is provided in \texttt{scripts/environment.yml} within the Zenodo deposit.

\subsection*{Computational Resources}

All computations were performed on the following infrastructure:

\begin{itemize}[leftmargin=*,itemsep=0.3em]
\item \textbf{CPU-based analyses:} High-performance computing cluster with Intel Xeon processors (16--32 cores per node, 64--128 GB RAM)
\item \textbf{GPU-accelerated training:} NVIDIA A4000 and RTX 4090 GPUs (16--24 GB VRAM) for GNN and Transformer model training
\item \textbf{Operating system:} Linux (Ubuntu 22.04 LTS)
\item \textbf{Job scheduling:} SLURM workload manager for distributed training and evaluation
\end{itemize}

\noindent Typical wall-clock times are documented in the Data Analysis Report, enabling researchers to estimate computational requirements for replication.

\subsection*{Data Access and Reuse}

\textbf{No restrictions:} All data and code are released under CC BY 4.0, permitting unrestricted reuse with attribution. Researchers may:

\begin{itemize}[leftmargin=*,itemsep=0.3em]
\item Reproduce all reported results using the archived splits and scripts
\item Extend the benchmark to additional models or molecular representations
\item Apply the trained models to new antimalarial screening campaigns
\item Adapt the evaluation framework to other therapeutic areas or chemical libraries
\end{itemize}

\subsection*{Data Format Specifications}

\begin{itemize}[leftmargin=*,itemsep=0.3em]
\item \textbf{Molecular structures:} CSV files with canonical SMILES strings (UTF-8 encoding)
\item \textbf{Activity labels:} Binary classification (0 = inactive, 1 = active) based on upstream data sources
\item \textbf{Chemical splits:} JSON files with integer indices mapping to panel rows
\item \textbf{Graph representations:} PyTorch Geometric \texttt{Data} objects serialized via \texttt{torch.save}
\item \textbf{Predictions:} CSV files with per-molecule probability scores and class labels
\item \textbf{Metrics:} JSON summaries with floating-point performance statistics
\item \textbf{Checksums:} SHA256 hashes in BSD-style format (\texttt{sha256sum} compatible)
\end{itemize}

\subsection*{Provenance and Quality Control}

The Data Analysis Report (\texttt{P5\_DATA\_ANALYSIS\_REPORT.md}) provides complete computational provenance:

\begin{itemize}[leftmargin=*,itemsep=0.3em]
\item SLURM job identifiers for all distributed computations
\item Execution timestamps and wall-clock durations
\item Quality-control gate results (fold completeness, finite metric verification)
\item Evidence status labels (\texttt{COMPUTED}, \texttt{COMPUTED\_SECONDARY}, \texttt{NOT\_COMPUTED})
\item Frozen input hashes preventing silent data substitution
\item Audit trails for all secondary analyses and post-hoc extensions
\end{itemize}

\noindent This granular provenance enables independent verification of computational integrity and facilitates debugging during replication attempts.

\subsection*{Contact Information}

For questions regarding data access, replication protocols, or model weights:

\begin{itemize}[leftmargin=*,itemsep=0.2em]
\item \textbf{Corresponding author:} Myke Vital Sao Temgoua (\texttt{myke-vital.sao@facsciences-uy1.cm})
%\item \textbf{GitHub issues:} \url{https://github.com/NanaEngo/Malaria_codesV2/issues}
\item \textbf{Zenodo contact:} Via Zenodo messaging system for deposit-specific inquiries
\end{itemize}

\subsection*{Long-Term Preservation}

The Zenodo deposit ensures long-term preservation through:

\begin{itemize}[leftmargin=*,itemsep=0.3em]
\item Permanent DOI assignment with guaranteed resolution
\item CERN-hosted infrastructure with institutional backing
\item Multiple geographic replicas for disaster resilience
\item Integration with DataCite metadata registry
\item Compliance with funding-agency open-data mandates
\end{itemize}

\subsection*{Related Datasets}

This study is part of an integrated antimalarial drug discovery program. Related datasets include:

\begin{itemize}[leftmargin=*,itemsep=0.3em]
\item \textbf{P1:} Chemical space exploration and multi-target docking (Zenodo DOI: \url{https://doi.org/10.5281/zenodo.22696778})
\item \textbf{P3:} Quantum-inspired molecular representations (Zenodo DOI: \url{https://doi.org/10.5281/zenodo.22768267})
\item \textbf{P2:} Molecular dynamics validation of polypharmacology (Zenodo DOI: \url{https://doi.org/10.5281/zenodo.22829031})
\end{itemize}

\noindent Cross-project integration is documented in the repository file \texttt{P1\_P7\_INTEGRATION\_ROADMAP.md}.

\subsection*{Compliance Statement}

This Data Availability Statement complies with:

\begin{itemize}[leftmargin=*,itemsep=0.2em]
\item Royal Society of Chemistry data-sharing policies
\item FAIR data principles (Findable, Accessible, Interoperable, Reusable)
\item ORCID author identification requirements (all co-authors have registered ORCIDs)
\item Open Science best practices for computational chemistry and machine learning
\end{itemize}

%\vspace{1em}
%\noindent \textbf{Version:} V2609 (September 2026)\\
%\textbf{Last updated:} \today

\end{document}
